\section{Introduction}
\label{sec:introduction}

The equality $3+17=7+13$ gives two decompositions of $20$ into odd
primes. Multiplying the corresponding monomials by $x_3$ gives a move
between decompositions of $23$ into three odd primes. A different move,
$3+7=5+5$, now gives
\[
 x_3^2x_{17}=x_3x_7x_{13}=x_5^2x_{13}.
\]
These equalities hold in the binary Goldbach ring over a field $\kk$
\[
 R=\kk[x_p:p\text{ an odd prime}],\qquad
 G=R/(x_px_q-x_{p'}x_{q'}\mid p+q=p'+q').
\]
Each move preserves the sum of the labels and the number of factors.
Accordingly, we give $x_p$ bidegree $(p,1)$ and call the two entries
the Goldbach degree and the monomial degree.
Goldbach's conjecture says that $G_{n,2}\ne0$ for every even $n\ge6$;
here we study the relations among these representations.
The monomials of $R_{n,d}$ correspond to unordered representations
of $n$ as a sum of $d$ odd primes. In $G_{n,d}$, two such monomials
give the same class precisely when a sequence of binary moves joins them.

Our main results concern interpolation and the growth of relations
among these representations. Equal pair sums force an affine
interpolation law for maps from the odd primes to any abelian group
(Theorem~\ref{main:prime-interpolation}), settling K\'atai and Phong's
conjecture on shifted prime arguments \cite[Conj.~6]{KP15Problems}.
Explicit Koszul cycles obtained by shifting primes span a positive
fraction of the linear syzygies in each sufficiently large admissible
Goldbach degree, and we determine the leading Betti asymptotic
(Theorem~\ref{main:weightwise}). Finally, the two natural finite
truncations have different leading constants for their total linear
strands (Theorem~\ref{main:cutoff-growth}).

The moves above already occur as edges of a strand of the minimal free
resolution of $G$ over $R$. In $F_{23,3}$ the vertices $x_3x_7x_{13}$,
$x_5^2x_{13}$, and $x_5x_7x_{11}$ form a triangle. Its oriented boundary
is a relation among the three edge generators, giving a minimal syzygy
(Example~\ref{ex:original23}). How does the number of such generators
grow with the Goldbach degree at a fixed monomial degree? Write
\[
 \beta_{i,n,d}=\dim_\kk\Tor_i^R(G,\kk)_{n,d},\qquad
 T_d(n)=\#\{(p_1,\ldots,p_d):p_j\text{ odd prime},\
                                  p_1+\cdots+p_d=n\}.
\]
A star in an internal degree means that we sum over that degree. Thus
$\beta_{i,*,d}=\sum_n\beta_{i,n,d}$ and
$\beta_i=\beta_{i,*,*}=\sum_{n,d}\beta_{i,n,d}$.
The count $T_d(n)$ includes repetitions. Write $\A$ for the odd primes
and $\A_t=\A\cap[1,t]$. We distinguish two growth questions. In the
full ring, we fix $d$ and let the Goldbach degree $n\to\infty$,
counting one bidegree at a time. For finite truncations on $\A_m$,
we fix $i$, sum over all Goldbach degrees in monomial degree $i+1$,
and let the cutoff $m\to\infty$.

For the first question, common shifts of primes give explicit Koszul
cycles in the linear strand. For example, adding two to each of $3,11,17$ gives
the primes $5,13,19$. In the Koszul complex on the variables this
common shift gives the cycle
\[
 x_5\varepsilon_{11}\varepsilon_{17}
 -x_{13}\varepsilon_3\varepsilon_{17}
 +x_{19}\varepsilon_3\varepsilon_{11}.
\]
Its differential vanishes by the binary relations in $G$. Up to sign,
this cycle represents the upper quadrilateral in the drawing of
$F_{33,3}$ in Example~\ref{ex:weight33}.
More generally, let $W^+_{n,d}$ be the span of the classes obtained
from $d$ distinct primes $p_1,\ldots,p_d$ and a shift $h>0$ for which
all $p_j+h$ are prime and $h+\sum_jp_j=n$. These classes have
homological degree $d-1$. The following estimates compare their span
with the entire linear strand.

\begin{maintheorem}[Growth in Goldbach degree and common shifts]
\label{main:weightwise}
Let $G$ be the binary Goldbach ring on all odd primes. Fix $d\ge3$.
Over every field,
\[
 \beta_{d-1,n,d}
 =\frac{d-1}{d!}\,T_d(n)+O_d\bigl((1+|\A_n|)^{d-2}\bigr).
\]
There are constants
$c_d,C_d>0$ independent of the field such that, for every sufficiently
large $n\equiv d\pmod2$,
\[
 c_d\,\frac{n^{d-1}}{(\log n)^d}
 \ \le\ \dim_\kk W^+_{n,d}
 \ \le\ \beta_{d-1,n,d}
 \ \le\ C_d\,\frac{n^{d-1}}{(\log n)^d}.
\]
\end{maintheorem}

In particular, for fixed $d\ge3$, as $n\to\infty$ through
$n\equiv d\pmod2$,
\[
 \beta_{d-1,n,d}\sim\frac{d-1}{d!}T_d(n)
 \sim\frac{\mathfrak S_d(n)}{d!(d-2)!}
             \frac{n^{d-1}}{(\log n)^d},
\]
where $\mathfrak S_d(n)$ is the prime singular series defined in
Corollary~\ref{cor:prime-betti-singular-series}. The estimates are
uniform over coefficient fields, and the translated cycles account
for a positive fraction of the linear strand in each such weight.
The coefficient $(d-1)/d!$ has a homological explanation.
For a partition into $d$ distinct primes there are $d$ ways to leave
one prime in the coefficient and put the other $d-1$ in the exterior
product. One incoming Koszul boundary relates these $d$ terms, leaving
$d-1$ dimensions before the outgoing differential is applied. Each
such partition has $d!$ orderings in $T_d(n)$. The error estimate
accounts for partitions with repeated primes and for the rank of the
outgoing differential. For $2\le k\le d$, if the defining moves replace
$k$ summands at a time, the same argument gives the
coefficient $\binom{d-1}{k-1}/d!$ in homological degree $d-k+1$
(Theorem~\ref{thm:fixed-weight-betti-euler}).

Theorem~\ref{thm:fixed-weight-betti-euler} proves the algebraic estimate
for arbitrary positive supports; Corollary~\ref{cor:prime-betti-singular-series}
combines it with known prime asymptotics. To account for a positive
fraction by explicit classes, Theorem~\ref{thm:prime-shift-proportion}
selects independent translated cycles by their exterior coordinates
and counts simultaneous prime completions using the linear forms
theorem of Tao and Ter\"av\"ainen \cite[Thm.~1.6]{TT}.
Related
translated cycles for lattice polygons appear in the work of
Castryck--Cools--Demeyer--Lemmens \cite[Secs.~4.1--4.2]{CCDL19}.
For a single shift, Proposition~\ref{thm:scroll-shift} in
Appendix~\ref{app:scroll-shift} uses the full scroll resolution to count
the resulting linear syzygies by Goldbach degree.

For the second growth question, $m$ bounds the available prime labels.
The variable cutoff $G^{<m}$ retains all binary relations on $\A_m$;
the relation-weight cutoff $G^{\le m}$ retains only those of common
weight at most $m$ on the same support. The superscripts distinguish
the ideals, while both support cutoffs are inclusive. Write
$\beta^{<m}$ and $\beta^{\le m}$ for their Betti numbers over
$R^{<m}=\kk[x_p:p\in\A_m]$.

\begin{maintheorem}[Linear-strand growth under truncation]
\label{main:cutoff-growth}
For every fixed $i\ge1$, as $m\to\infty$,
\[
 \beta^{<m}_{i,*,i+1}
       \sim\frac{i}{(i+1)!}|\A_m|^{i+1},\qquad
 \beta^{\le m}_{i,*,i+1}
       \sim\frac{1-2^{-i}}{(i+1)!}|\A_m|^{i+1},
\]
uniformly over coefficient fields.
\end{maintheorem}

Here $|\A_m|\sim m/\log m$, and each total sums over all Goldbach
degrees $n$ while keeping monomial degree $i+1$ fixed. For cubic
second syzygies the leading constants are $1/3$ and $1/8$.
Theorem~\ref{grow:both-cutoffs} proves these formulas and their
analogues for every fixed arity. Their proof counts across the finite
support and accounts for which relation weights are retained.
In contrast, Theorem~\ref{main:weightwise} concerns one bidegree of
the full ring. Stabilization identifies each fixed bidegree with its
finite-cutoff value once $m$ is large enough
(Proposition~\ref{setup:stabilization}).

The same binary relations also determine which functions on the odd
primes preserve equal pair sums.

\begin{maintheorem}[Interpolation on the odd primes]\label{main:prime-interpolation}
Let $\Podd$ be the set of odd primes, let $H$ be an abelian group,
and let $g:\Podd\to H$ satisfy
$g(p)+g(q)=g(p')+g(q')$ whenever $p+q=p'+q'$. Then for every odd prime $p$,
\[
 g(p)=g(3)+\frac{p-3}{2}\bigl(g(5)-g(3)\bigr).
\]
\end{maintheorem}

We prove this by computing the universal Freiman group generated by the
odd primes modulo their equal-pair-sum relations
(Theorem~\ref{depth:prime-pair-group}). It is free of rank two.
The proof adjoins auxiliary prime labels to realize a prescribed
equal-sum identity by binary moves, and then cancels those auxiliary
classes in the group. This converts the problem into a system of
linear forms to which the prime-counting theorem applies.
To our knowledge, this explicit integral identification has not
previously been recorded. It determines the group itself, including
the absence of torsion, and gives interpolation in every abelian target.

The complex-valued case proves the odd-prime affine-classification
assertion in Conjecture~1 of K\'atai--Phong \cite[p.~469]{KP15}, who prove affine classification
for functions preserving sums of at least three primes and propose
a binary analogue. Park's survey \cite[p.~6, after Thm.~3.7]{Park25}
still records the binary analogue as a conjecture.
Theorem~\ref{main:prime-interpolation} extends this odd-prime
classification to arbitrary abelian targets. The restriction to odd
primes matters: the printed all-primes formulation leaves an
independent value at $2$ under its pair-sum hypothesis, as explained
after Theorem~\ref{depth:prime-pair-group}.
Applying our interpolation formula to shifted arguments also proves
the distinct Conjecture~6 of K\'atai--Phong
\cite[p.~235]{KP15Problems}, as shown in Corollary~\ref{depth:KP-interpolation}.

Translation classes also occur in arbitrarily high homological
degrees. Unbounded multiplicities of positive prime differences imply
$\pd_R G=\infty$ (Corollary~\ref{cor:density}); arbitrarily long prime
progressions give another route through Green--Tao \cite{GreenTao08}.

The connection between additive structure and commutative algebra has
been developed through Hilbert functions of iterated sumsets
\cite{EliahouMazumdar22}, projective monomial curves \cite{Elias22},
and the interpretation of Freiman's doubling inequality through
monomial ideals and their fiber cones \cite{HerzogHibiZhu}.
The connectivity of monomials under chosen moves also fits the
binomial-ideal framework of Diaconis--Eisenbud--Sturmfels \cite{DES98}.
To distinguish our ring from the full toric algebra, let
$J=\ker(R\to\kk[u,v])$, where $x_p\mapsto u^pv$.
The full toric quotient $R/J$ identifies all monomials of the same
bidegree, whereas the binary Goldbach ring is
$G=R/\langle J_{*,2}\rangle$, with $J_{*,2}=\bigoplus_n J_{n,2}$.
Thus $G$ imposes the quadratic equal-sum relations and their ideal
consequences, retaining the components under pair replacements.
The natural surjection $G\to R/J$ is an isomorphism precisely when
every nonempty binary-move fiber is connected.
For the full odd-prime support, we ask whether this is always so, or
equivalently whether $G$ is a domain
(Question~\ref{question:full-support-connectivity}).
For general supports connectivity can fail. We study how the fiber
components and the prime representation counts $T_d(n)$ enter the
minimal resolution.

The group computation also controls the finite variable-cutoff
resolutions. For all sufficiently large $m$, the rank-two group gives
\[
 \depth G^{<m}\le2,\qquad
 |\A_m|-2\le\pd_{R^{<m}}G^{<m}\le|\A_m|.
\]
An exhaustive prime-pair verification also gives
$\depth G^{<m}\le2$ at every integer cutoff
$3\le m\le20\,000\,000$, over every field
(Corollary~\ref{depth:verified-range}). This finite statement uses
triangular relations and localization, without computing the minimal
resolution; it does not determine the exact depth.
Theorem~\ref{depth:relation-cutoff-pd} gives an asymptotically maximal
projective dimension for the relation-weight cutoff as well.
Group completion identifies monomials of the same bidegree after
inverting the variables; it does not itself allow cancellation in $G$
or prove connectivity of every binary-move fiber.

The Betti estimates also have a Hilbert-series consequence. Put
\[
 \gamma_{n,d}=\sum_i(-1)^i\beta_{i,n,d},\qquad
 E_G(u,v)=\sum_{n,d}\gamma_{n,d}u^nv^d.
\]
If $H_G$ and $H_R$ denote the bigraded Hilbert series, the Euler
identity gives $E_G=H_G/H_R$. Corollary~\ref{cor:hilbert-euler-comparison}
shows that $(-1)^{d-1}\gamma_{n,d}$ differs from the dominant binary
Betti number by a lower-order term for fixed $d\ge3$ as
$n\to\infty$ through $n\equiv d\pmod2$. Thus the prime
representation asymptotic also governs coefficients determined by the
Hilbert series. The component basis makes $H_G$, and hence $E_G$,
independent of the field, although individual Betti numbers can vary.

The companion paper \cite{CompanionTopology} studies residue-field
Betti numbers, integral torsion, and higher operations in these rings.
Here the emphasis is on the arithmetic estimates and their consequences
for interpolation, depth, and projective dimension.
Throughout, $\kk$ denotes a field of characteristic $\ell$, where
$\ell=0$ or $\ell$ is prime. We reserve $k$ for the arity of a
defining additive relation. The variable $n$ is the Goldbach degree, $d$ the monomial degree, $i$ the
homological degree, and $m$ the truncation parameter.

Sections~\ref{sec:setup}--\ref{sec:koszul} establish the algebraic and
geometric framework. Section~\ref{sec:betti-euler-shifts} proves the
fixed-degree Betti estimates and the positive proportion of shifted
cycles as $n\to\infty$, then records their Hilbert-series consequence.
Section~\ref{sec:depth} treats interpolation and depth, and
Section~\ref{sec:growth} proves the truncation estimates as $m\to\infty$.
Section~\ref{sec:questions} concludes with further questions.
Appendix~\ref{app:scroll-shift} gives the scroll bound for a fixed
shift, and Appendix~\ref{app:saturation} studies saturation, quotients
by a prime variable, and dual Koszul cocycles.
